% SOURCE_AUTHORITY: sga5_fr_workpass.tex, lines 13614--13640 (LF-normalized unit).
% SOURCE_SHA256: 791F4EFFC5E02832D5D77ED03518C8156D6F07E4C8238B03545DB93D883FBB28
Se dotarmos o espaço vetorial sobre $k$
\[
\Hom_k(\fm/\fm^2 \; , \; \fm^r/\fm^{r+1})
\]
da estrutura habitual de $G$-módulo, seu elemento $u(f'_R)$ pertence ao subespaço
\[
\Hom_k(\fm/\fm^2 \; , \; \fm^r/\fm^{r+1})^G
\]
dos invariantes sob $G$.

Por outro lado, é clássico que o espaço vetorial $\fm^i/\fm^{i+1}$ se identifica canonicamente com a $i$-ésima potência tensorial $\bigotimes_k^i (\fm/\fm^2)$ e, por conseguinte, o espaço vetorial $\Hom_k(\fm/\fm^2 \; , \; \fm^r/\fm^{r+1})$ se identifica mediante o homomorfismo de «contração com a $(r-1)$-ésima potência tensorial $\bigotimes_k^{r-1}(\fm/\fm^2)$»:
\[
\gamma : \Hom_k(\fm/\fm^2 \; , \; \fm^r/\fm^{r+1}) \longrightarrow \bigotimes_k^{r-1}(\fm/\fm^2) \quad .
\]

Por meio do isomorfismo $\gamma$, a estrutura de $G$-módulo de $\Hom_k(\fm/\fm^2 \; , \; \fm^r/\fm^{r+1})$ é transportada para a estrutura de $G$-módulo de $\bigotimes_k^{r-1}(\fm/\fm^2)$ induzida pela estrutura de $G$-módulo de $\fm/\fm^2$:
\[
g(\xi_1 \otimes \xi_2 \otimes \cdots \otimes \xi_{r-1}) = g\xi_1 \otimes \cdots \otimes g\xi_{r-1} \quad .
\]

Denotaremos por $\mu(f'_R)$ a imagem de $u(f'_R)$ pelo isomorfismo $\gamma$.

Tem-se, portanto:
\begin{enumerate}
\item[I)] $0 \neq \mu(f'_R) \in \bigotimes_k^{r-1}(\fm/\fm^2)$,
\item[II)] $\mu(f'_R)$ é invariante sob $G$.
\end{enumerate}
